\begin{figure}[t!]
\centering
% Certifique-se de incluir no preâmbulo:
% \usetikzlibrary{positioning, calc, arrows.meta}
\begin{tikzpicture}[
    scale=1.0, 
    thick, >=stealth,
    % Caixa delimitadora externa da célula
    cell_box/.style={draw, rectangle, rounded corners=10pt, black!80, loosely dashed, line width=1.2pt, minimum width=7.2cm, minimum height=4.2cm},
    % Blocos de ativação
    layer/.style={draw, rectangle, font=\scriptsize\bfseries, minimum width=1.0cm, minimum height=0.5cm, fill=wp-accent, text=black},
    % Operadores matemáticos circulares
    op/.style={draw, circle, minimum size=0.5cm, inner sep=0pt, font=\small, fill=wp-primary!50, draw=black},
    % Estilo de texto interno para os rótulos posicionados sobre as linhas
    gate_lbl/.style={font=\small\bfseries, text=black, align=center, fill=white, inner sep=2pt},
    io/.style={font=\small, text=black}
]

    % 1. Renderiza a caixa externa principal (Referência espacial)
    \node[cell_box] (B) at (4.0, 1.8) {};

    % 2. Nós de entrada externos e rótulos (O fluxo agora se conecta diretamente a eles)
    \node[io, left=0.8cm of B.155, align=left] (c_prev_lbl) {Memory cell\\internal state\\$\mathbf{C}_{t-1}$};
    \node[io, left=0.8cm of B.190, align=left] (h_prev_lbl) {Hidden state\\$\mathbf{H}_{t-1}$};
    \node[io, below=0.6cm of B.225] (x_lbl) {Input $\mathbf{X}_t$};
    
    % Saída à direita
    \node[io, right=0.8cm of B.25] (c_next_lbl) {$\mathbf{C}_t$};

    % 3. Alinhamento dos Operadores na Esteira Superior
    \node[op] (f_mult) at ($(B.155)!0.22!(B.25)$) {$\odot$};
    \node[op] (i_add)  at ($(B.155)!0.65!(B.25)$) {$+$};

    % 4. Posicionamento Estratégico dos Blocos de Ativação
    \node[layer] (f_gate) at ($(f_mult)-(0,1.8)$) {$\sigma$};
    \node[op] (input_mult) at ($(i_add)-(0,0.9)$) {$\odot$};
    \node[layer] (c_gate) at ($(input_mult)-(0,0.9)$) {$\tanh$};
    \node[layer] (i_gate) at ($(c_gate)-(1.6,0)$) {$\sigma$};
    \node[layer] (o_gate) at ($(B.155)!0.85!(B.25) - (0,1.8)$) {$\sigma$};

    % ================= CONEXÕES COM RÓTULOS NO MEIO DAS RETAS =================

    % Fluxo Principal Superior: Cell State (Setas saem e entram DIRETAMENTE nas letras)
    \draw[->] (c_prev_lbl.east) -- (f_mult);
    \draw[->] (f_mult) -- (i_add);
    \draw[->] (i_add) -- (c_next_lbl.west);

    % Linha de Barramento Inferior (Hidden State)
    \coordinate (h_bus) at ($(B.190)+(0.3,0)$);
    \draw[-] (h_prev_lbl.east) -- (h_bus);
    \draw[-] (x_lbl.north) |- (h_bus);

    % Distribuição paralela para a base de cada camada
    \foreach \g in {f_gate, i_gate, c_gate, o_gate} {
        \draw[->] (h_bus) -| (\g.south);
    }

    % Forget gate à esquerda da sua linha vertical
    \draw[->] (f_gate) -- (f_mult) node[midway, left=2pt, font=\tiny\bfseries, align=center] {Forget\\gate\\$\mathbf{F}_t$};
    
    % Input node à direita da sua linha vertical
    \draw[->] (c_gate) -- (input_mult) node[midway, right=2pt, font=\tiny\bfseries, align=center] {Input\\node\\$\tilde{\mathbf{C}}_t$};
    
    % Input gate no meio da linha ortogonal que dobra em L
    \draw[->] (i_gate.north) |- (input_mult.west) node[midway, above, font=\tiny\bfseries, yshift=-1pt] {Input gate $\mathbf{I}_t$};
    
    % Output gate no meio da linha vertical que sobe livremente
    \draw[->] (o_gate.north) -- ($(o_gate.north)+(0,1.0)$) node[midway, fill=white, inner sep=1pt, font=\tiny\bfseries] {Output gate $\mathbf{O}_t$};

\end{tikzpicture}
\caption{Corrected schematic of an LSTM cell. The memory flow arrows connect directly to the text variables, and all gate labels are mathematically centered along their internal routing lines.}
\label{fig:lstm_perfect_routing}
\end{figure}